
\subsubsection{2.1 LLM Self-Repair}

Iterative self-repair feeds execution feedback back to the model to enable code improvement. Olausson et al. (2023) demonstrated 12-17\% relative improvement on HumanEval. Arimbur (2026) showed that most gains concentrate in the first 2-3 repair iterations.

Several frameworks operationalize self-repair. CodeRL integrates reinforcement learning with execution feedback. ThinkRepair uses self-directed debugging on Defects4J. CodeCoR (2025) achieves 77.13\% pass@1 on HumanEval/MBPP through agent collaboration.

\subsubsection{2.2 Static Analysis Feedback}

Static analysis provides signals complementary to execution feedback. Jain et al. (2024) showed that LLM-assisted code cleaning improves code quality without runtime signals. Blyth et al. (2025) demonstrated that Pylint feedback reduces security issues from >40\% to 13\% and reliability warnings from >50\% to 11\%.

AutoSafeCoder combines static analysis with fuzz testing, achieving 13\% vulnerability reduction. Dolcetti et al. (2024) integrated testing and static analysis feedback for safety improvements.

Studies comparing static and execution feedback face a confound: they necessarily vary information volume alongside feedback type. Cascaded approaches (static + execution) provide more information than single-type approaches.

\subsubsection{2.3 Feedback Presentation Effects}

Research on prompt engineering and positional effects suggests LLMs are sensitive to presentation order. Recency effects in transformer attention mean later tokens may receive disproportionate weight. FeedbackEval benchmarked feedback-driven code repair, finding mixed feedback yields 63.6\% repair success, though ordering effects were not systematically tested.

